\ref{ms:lemma-12-cap} 的局部延拓界为 $O_\varepsilon(R^\varepsilon)$，对任意 $\varepsilon>0$ 都成立。同样的估计在一个固定的较大参数球上也成立。实际上，作变换从 $v$ 到 $Lv$，其中 $L>1$ 固定，会使物理坐标按 $(x,t)\mapsto (Lx,L^2t)$ 变换，半径至多增大一个固定倍数。在较大面片上取光滑截断，使它在 $B_d$ 上等于一，再乘以参数测度 $\mathrm dv$，可得光滑紧支撑曲面密度 $\mu$，满足
\[
|\widehat\mu(x,t)|\lesssim (1+|(x,t)|)^{-d/2}.
\]
当 $|x|$ 远大于 $|t|$ 时，用分部积分。当上述情形不成立且 $|(x,t)|$ 较大时，相位关于 $v$ 的 Hessian 为 $2t I_d$，其中 $|t|\gtrsim |(x,t)|$，由驻相法得到该界。

由对偶性，局部估计给出
\[
\|\widehat F(v,|v|^2)\|_{L^{p_c'}(\mu)}
 \lesssim_\varepsilon R^\varepsilon\|F\|_{L^{p_c'}(\mathbb R^{d+1})}
\]
对光滑函数 $F$ 均成立，只要其支撑包含在任意一个半径为 $R\geq1$ 的球内。为了看出截断不改变局部假设，将输入与截断函数相乘，再应用较大面片的延拓算子，即得到以 $\mu$ 为测度的延拓。截断函数有界，因此输入的 $L^{p_c}(\mu)$ 范数控制乘积相应的 $L^{p_c}(\mathrm dv)$ 范数。应用 Tao 的 epsilon 消去定理 \cite{r-75}, Theorem~1.2，采用 Kim \cite{r-46}, Theorem~5.3 针对紧超曲面的表述。对具有幂次 Fourier 衰减的紧光滑超曲面测度，若在每个球上都有 $1<u\leq v\leq 2$ 的局部限制估计，允许任意小的幂次损失，则可推出全局限制从 $L^w$ 到 $L^v(\mu)$，其中 $1\leq w<u$。这里 $u=v=p_c'\in(1,2)$。因此，使用对偶并将 $B_d$ 上的密度以零延拓，得到 $E:L^{p_c}(B_d)\to L^q(\mathbb R^{d+1})$，对每个 $p_c<q<\infty$ 成立。由 $B_d$ 的测度有限，此范围内有 $E:L^q\to L^q$。

若 $a\geq b$，用 $q=b$ 和 $B_d$ 的测度有限。若 $a<b$，令 $q=1+b/a'$。假设给出 $q>1+(d+2)/d=p_c$，而 $a<b$ 给出 $q<b$。对 $E:L^1\to L^\infty$ 和 $E:L^q\to L^q$ 用 $\theta=q/b$ 插值。于是 $1/b=\theta/q$ 和 $1/a=1-\theta+\theta/q$。这与 \ref{ms:cor-9-3} 对球面所作的插值计算相同。
